%% SYNTHETIC TEST UNIT generated by writing/tools/make_outline_book.py from FINAL_OUTLINE_v2 -- not book text
\chapter{Drell–Yan 过程与强子—强子碰撞}\label{ch:dy}
本单元为测试用的合成文本（由 FINAL\_OUTLINE\_v2 自动生成），陈述内容从略。

\section{第 1 节（测试）}\label{sec:dy-born}

\begin{definition}[DY 运动学]\label{def:dy-kin}
（测试用合成陈述：内容从略。）
\uses{def:prelim-lightfront}
\end{definition}

\begin{proposition}[Born 截面 \tagR]\label{prop:dy-born}
（$R_{\mathrm{PT}}$：在重整化微扰论中对红外正规化的部分子态严格。）
（测试用合成陈述：内容从略。）
\uses{def:dy-kin, def:pdf-quark}
\end{proposition}

\begin{proposition}[DY 的领头区域 \tagA]\label{prop:dy-regions}
（测试用合成陈述：内容从略。）
\uses{thm:ir-leading-regions, def:dy-kin}
\end{proposition}
\begin{proofsketch}测试。\end{proofsketch}
\paragraph{地位与缺口}测试：本陈述的地位由\cref{asm:appN-envelope}、\cref{asm:appN-separation}决定。

\begin{proposition}[单个割图中的 Glauber 钳夹 \tagR]\label{prop:dy-glauber-pinch}
（$R_{\mathrm{PT}}$：在重整化微扰论中对红外正规化的部分子态严格。）
（测试用合成陈述：内容从略。）
\uses{prop:ir-glauber-trapping, def:ir-glauber}
\end{proposition}

\section{第 2 节（测试）}\label{sec:dy-regions}

\begin{lemma}[光锥序中间态的取向与 $k^-$ 极点 \tagR]\label{lem:dy-poles-orientation}
（$R_{\mathrm{PT}}$：在重整化微扰论中对红外正规化的部分子态严格。）
（测试用合成陈述：内容从略。）
\uses{lem:appO-lcopt, lem:coll-final-poles, prop:dy-glauber-pinch, def:appO-lcopt, def:appO-orientation, lem:appO-orientation}
\end{lemma}

\begin{remark}[每个喷注一个光锥时间]\label{rem:dy-one-time-per-jet}
（测试用合成陈述：内容从略。）
% uses: none
\end{remark}

\begin{lemma}[旁观者末态相互作用的幺正性相消 \tagA]\label{lem:dy-fsi}
（测试用合成陈述：内容从略。）
\uses{prop:method-fsi, lem:appO-unitarity, asm:appO-lcopt}
\end{lemma}
\begin{proofsketch}测试。\end{proofsketch}
\paragraph{地位与缺口}测试：本陈述的地位由\cref{asm:appO-lcopt}决定。

\begin{example}[单个割图给出未来指向的 Wilson 线 \tagR]\label{ex:dy-future-vs-past}
（$R_{\mathrm{def}}$：由算符定义直接推出。）
（测试用合成陈述：内容从略。）
\uses{prop:prelim-wilson-rules}
\end{example}

\section{第 3 节（测试）}\label{sec:dy-orientation}

\begin{proposition}[移除末态极点后的围道变形 \tagA]\label{prop:dy-deformation}
（测试用合成陈述：内容从略。）
\uses{prop:method-glauber-deformation, lem:dy-fsi, lem:dy-poles-orientation, lem:appO-unitarity, lem:appO-equal-time, lem:appO-joint-localization, prop:appO-one-sided, prop:appO-target, prop:appO-eikonal, prop:appO-past-wilson, asm:appO-mixing, asm:appO-deformation, asm:appO-envelope, asm:appO-rewritten-cs}
\end{proposition}
\begin{proofsketch}测试。\end{proofsketch}
\paragraph{地位与缺口}测试：本陈述的地位由\cref{asm:appO-mixing}、\cref{asm:appO-deformation}、\cref{asm:appO-envelope}、\cref{asm:appO-rewritten-cs}决定。

\begin{lemma}[软因子为 1 \tagA]\label{lem:dy-soft}
（测试用合成陈述：内容从略。）
\uses{prop:dy-deformation, lem:method-soft-unitarity, prop:appO-past-wilson}
\end{lemma}
\begin{proofsketch}测试。\end{proofsketch}
\paragraph{地位与缺口}测试：本陈述的地位由\cref{prop:dy-deformation}决定。

\begin{proposition}[DY 与 DIS 的共线 PDF 是同一对象 \tagR]\label{prop:dy-universality}
（$R_{\mathrm{def}}$：由算符定义直接推出。）
（测试用合成陈述：内容从略。）
\uses{prop:prelim-wilson-covariance, def:pdf-quark}
\end{proposition}

\begin{proposition}[阿贝尔模型中旁观者 Glauber 相位的相消 \tagP]\label{prop:dy-abelian-model}
（测试用合成陈述：内容从略。）
\uses{prop:method-glauber-oneloop, lem:prelim-eikonal-identity}
\end{proposition}
\paragraph{地位与缺口}测试：见文献。

\section{第 4 节（测试）}\label{sec:dy-glauber}

\begin{example}[旁观者模型的光锥序表示]\label{ex:dy-spectator-lcopt}
（测试用合成陈述：内容从略。）
\uses{lem:appO-lcopt, lem:dy-poles-orientation, prop:dy-abelian-model}
\end{example}

\begin{example}[阿贝尔 Glauber 模型的数值演示]\label{ex:dy-glauber-model}
（测试用合成陈述：内容从略。）
\uses{prop:dy-abelian-model}
\end{example}

\begin{proposition}[两条程函线之间的双胶子交换：$O(g^4)$ 幺正性与 Regge 相位 \tagR]\label{prop:dy-eikonal-twogluon}
（$R_{\mathrm{PT}}$：在重整化微扰论中对红外正规化的部分子态严格。）
（测试用合成陈述：内容从略。）
\uses{prop:prelim-wilson-rules, lem:qcd-colour-identities, lem:appO-glauber-operator}
\end{proposition}

\begin{remark}[共同时序是物理结果]\label{rem:dy-common-ordering}
（测试用合成陈述：内容从略。）
% uses: none
\end{remark}

\section{第 5 节（测试）}\label{sec:dy-model}

\begin{theorem}[单举 DY 的共线因子化 \tagA]\label{thm:dy-factorization}
（测试用合成陈述：内容从略。）
非微扰地位：\cref{asm:intro-bridge}。
\uses{thm:method-template, prop:dy-regions, lem:dy-fsi, prop:dy-deformation, lem:dy-soft, prop:dy-universality, thm:appO-main, cor:appO-interface, prop:appN-dy, def:pdf-quark, def:pdf-gluon, asm:intro-bridge, asm:appO-lcopt, asm:appO-envelope, asm:appO-mixing, asm:appO-deformation, asm:appO-rewritten-cs, asm:appN-collins, asm:appN-cutsum}
\end{theorem}
\begin{proofsketch}测试。\end{proofsketch}
\paragraph{地位与缺口}测试：本陈述的地位由\cref{thm:appO-main}、\cref{cor:appO-interface}、\cref{prop:appN-dy}、\cref{asm:appO-lcopt}、\cref{asm:appO-envelope}、\cref{asm:appO-mixing}、\cref{prop:appO-twoexch-spectator}、\cref{asm:appO-deformation}、\cref{asm:appO-rewritten-cs}、\cref{asm:appN}、\cref{asm:intro-bridge}、\cref{asm:appN-brst-nonlocal}、\cref{thm:method-template}、\cref{prop:dy-eikonal-twogluon}、\cref{rem:appO-literature-comparison}决定。

\begin{remark}[固定阶部分子截面不检验旁观者 Glauber 相消]\label{rem:dy-nnlo-not-glauber}
（测试用合成陈述：内容从略。）
% uses: none
\end{remark}

\begin{remark}[Doria–Frenkel–Taylor 非相消的前提]\label{rem:dy-dft-premises}
（测试用合成陈述：内容从略。）
% uses: none
\end{remark}

\begin{proposition}[$O(\alpha_s)$ 的 $\overline{\rm MS}$ 系数函数 \tagR]\label{prop:dy-nlo}
（$R_{\mathrm{PT}}$：在重整化微扰论中对红外正规化的部分子态严格。）
（测试用合成陈述：内容从略。）
（本命题只涉及部分子态。）
\uses{def:dy-kin, prop:sudakov-oneloop, lem:prelim-plus-expansion, prop:pdf-partonic-dr, lem:appC-phase-space}
\end{proposition}

\section{第 6 节（测试）}\label{sec:dy-theorem}

\begin{example}[$O(\alpha_s)$ 系数函数的完整推导]\label{ex:dy-nlo}
（测试用合成陈述：内容从略。）
\uses{prop:dy-nlo}
\end{example}

\begin{proposition}[$\mu_F$ 无关性 \tagR]\label{prop:dy-scale}
（$R_{\mathrm{PT}}$：在重整化微扰论中对红外正规化的部分子态严格。）
（测试用合成陈述：内容从略。）
\uses{prop:dy-nlo, thm:pdf-dglap}
\end{proposition}

\begin{remark}[DIS 方案中的 DY 与 PDF 的普适性]\label{rem:dy-dis-scheme}
（测试用合成陈述：内容从略。）
% uses: none
\end{remark}

\begin{example}[大 $N$ 极限]\label{ex:dy-threshold}
（测试用合成陈述：内容从略。）
\uses{prop:dy-nlo, lem:prelim-mellin}
\end{example}

\section{第 7 节（测试）}\label{sec:dy-nlo}

\begin{theorem}[色单态末态 \tagA]\label{thm:dy-colour-singlet}
（测试用合成陈述：内容从略。）
非微扰地位：\cref{asm:intro-bridge}。
\uses{thm:dy-factorization, asm:intro-bridge}
\end{theorem}
\begin{proofsketch}测试。\end{proofsketch}
\paragraph{地位与缺口}测试：本陈述的地位由\cref{thm:dy-factorization}决定。

\begin{theorem}[极化 Drell–Yan \tagA]\label{thm:dy-polarized}
（测试用合成陈述：内容从略。）
非微扰地位：\cref{asm:intro-bridge}。
\uses{thm:dy-factorization, def:pdf-polarized, asm:intro-bridge}
\end{theorem}
\begin{proofsketch}测试。\end{proofsketch}
\paragraph{地位与缺口}测试：本陈述的地位由\cref{thm:dy-factorization}决定。

\begin{conjecture}[强子对撞中有色末态的共线因子化 \tagC]\label{conj:dy-coloured-final}
（测试用合成陈述：内容从略。）
\uses{thm:dy-factorization}
\end{conjecture}
\paragraph{地位与缺口}测试：本陈述的地位由\cref{prop:break-measurement}决定。

\begin{conjecture}[极化强子对撞中的 $A_{LL}$（喷注、$\pi^0$） \tagC]\label{conj:dy-polarized-jets}
（测试用合成陈述：内容从略。）
\uses{conj:dy-coloured-final, thm:dy-polarized}
\end{conjecture}
\paragraph{地位与缺口}测试：本陈述的地位由\cref{conj:dy-coloured-final}决定。

\section{第 8 节（测试）}\label{sec:dy-other}

\begin{theorem}[$1/Q^2$ 处的广义因子化 \tagA]\label{thm:dy-higher-twist}
（测试用合成陈述：内容从略。）
非微扰地位：\cref{asm:intro-bridge}。
\uses{thm:dy-factorization, asm:intro-bridge}
\end{theorem}
\begin{proofsketch}测试。\end{proofsketch}
\paragraph{地位与缺口}测试：见文献。

\begin{proposition}[更高幂次的因子化失效 \tagV]\label{prop:dy-twist6-failure}
（测试用合成陈述：内容从略。）
\uses{thm:dy-higher-twist}
\end{proposition}
\paragraph{地位与缺口}测试：见文献。
